\normalsize{El apartado ``Requisitos de las acreditaciones internacionales EUR-ACE y ABET'' de la
Normativa de TFT de la ETSIT-UPM establece que ``\textit{La memoria del TFT del GITST, GIB y MUIT,
y en general la de aquellas titulaciones que hayan obtenido o para las que se desee solicitar una
acreditación internacional EUR-ACE o ABET, debe mostrar conciencia de la responsabilidad de la
aplicación práctica de la ingeniería, el impacto social y ambiental, el compromiso con la ética
profesional, la responsabilidad y las normas de la aplicación práctica de la ingeniería, así como
sobre las prácticas de gestión de proyectos, gestión, y control de riesgos, entendiendo sus
limitaciones}''.

Este anexo obligatorio del TFT tendrá un carácter sintético con los siguientes apartados:}
\vspace{0.1cm}

\Large{\textsc{\textcolor{nar}{A1. Introducción}}} \par
\normalsize{Este trabajo desarrolla Voctomix~2.0, un sistema de producción remota de vídeo en directo
basado íntegramente en software libre (GStreamer, Python, Docker, RabbitMQ). El proyecto responde a
la necesidad de democratizar el acceso a herramientas de producción audiovisual profesional,
eliminando la dependencia de hardware dedicado de alto coste. Los principales asuntos identificados
son: impacto en la accesibilidad económica y social a la producción audiovisual, responsabilidad
ética en el uso de imágenes y datos durante retransmisiones en directo, y reducción del impacto
ambiental mediante la sustitución de hardware especializado por software sobre equipos
de propósito general.}
\vspace{0.1cm}

\Large{\textsc{\textcolor{nar}{A2. Descripción de impactos relevantes relacionados con el proyecto}}} \par
\vspace{0.3cm}
\normalsize{Los impactos más relevantes identificados son:
\begin{itemize}
  \item \textbf{Social:} Voctomix~2.0 permite a comunidades académicas, asociaciones sin ánimo
        de lucro y grupos de investigación como CyberNemo (UPM) realizar producciones audiovisuales
        de calidad sin presupuestos elevados. El sistema es completamente auditable y redistribuible
        bajo licencia GPL, lo que fomenta la transparencia y la colaboración.
  \item \textbf{Económico:} El sistema reemplaza equipos hardware cuyo coste puede superar los
        5.000~\euro\ (mezcladores Blackmagic, matrices de vídeo, sistemas de grabación) por software
        ejecutado en hardware convencional, con un coste de despliegue prácticamente nulo.
  \item \textbf{Ambiental:} La reducción de hardware especializado disminuye la fabricación de
        equipos dedicados y prolonga la vida útil de los equipos existentes. El despliegue en
        contenedores Docker facilita la reutilización de infraestructura compartida.
  \item \textbf{Ético-legal:} El uso del sistema para retransmisiones en directo implica la
        responsabilidad de respetar los derechos de imagen de las personas grabadas y los derechos
        de autor sobre el contenido audiovisual emitido.
\end{itemize}}
\vspace{0.1cm}

\Large{\textsc{\textcolor{nar}{A3. Análisis detallado de alguno de los impactos}}} \par
\vspace{0.3cm}
\normalsize{El impacto social más significativo es la \textbf{democratización del acceso a la
producción audiovisual profesional}. Hasta la fecha, una instalación multicámara profesional
requería inversiones de varios miles de euros en hardware dedicado. Voctomix~2.0, al ser
desplegable con un único comando (\texttt{docker compose up}), elimina esta barrera económica
y técnica, permitiendo que cualquier organización con un ordenador convencional pueda gestionar
producciones en directo de calidad. La arquitectura basada en contenedores garantiza además la
reproducibilidad del entorno, reduciendo la dependencia de configuraciones manuales complejas
y la necesidad de personal técnico altamente especializado.}

\Large{\textsc{\textcolor{nar}{A4. Conclusiones}}} \par
\vspace{0.3cm}
\normalsize{Voctomix~2.0 es un proyecto con un balance ético, social, económico y ambiental
positivo. El uso de software libre como base tecnológica garantiza la transparencia y la
continuidad del proyecto más allá del trabajo individual. La eliminación de dependencias de
hardware especializado reduce tanto el coste económico de acceso a la producción audiovisual
como el impacto ambiental asociado a la fabricación de equipos dedicados. Los criterios de
sostenibilidad adoptados, especialmente la contenerización Docker y la licencia GPL, aportan
valor añadido al proyecto al facilitar su adopción, mantenimiento y extensión por parte de
la comunidad.}
